% GENERATED FILE. Edit canonical lessons, then run npm run generate:books.
% canonical-source-hash: fnv1a64:e110b89a10c91953
% canonical-lessons: AR-C443-bawwab, AR-C443-mulaa, AR-C443-nazafa, AR-C443-atalla, AR-C443-khazna

\chapter{Bed and view}
\label{ch:ar-bed-and-view}
\hlchaptermodality{443}

\begin{chapteropening}
\textbf{By the end of this chapter:} \emph{I can say a doorman, a bed sheet, cleanliness, to overlook and a strongbox.}

Complete the last lesson of chapter 443: I can say a doorman, a bed sheet, cleanliness, to overlook and a strongbox.
\end{chapteropening}

\section[\texorpdfstring{\ar{بواب} (bawwāb)}{بواب (bawwāb)}]{\ar{بواب} (bawwāb) — a doorman; a caretaker}

\label{lesson:AR-C443-bawwab}



\begin{quote}
Before the new one: say the Arabic for to set alight, then the Arabic for pine trees.
\end{quote}

\subsection*{You'll want to know: \ar{بواب}}
\textbf{\ar{بواب}} — \emph{bawwāb} — \textquotedblleft{}a doorman; a caretaker\textquotedblright{}.

\begin{quote}
\ar{حمل} \ar{البواب} \ar{أمتعتنا} \ar{إلى} \ar{المصعد}.
\end{quote}

\begin{quote}
\emph{ḥamala al-bawwāb amtiʿatanā ilā al-miṣʿad.}
\end{quote}

\begin{quote}
The doorman carried our luggage to the lift.
\end{quote}

\begin{cousinweb}
\textbf{\ar{بواب}} (\emph{bawwāb}) is built on the native Arabic root \emph{b-w-b}, \textquotedblleft{}a door\textquotedblright{} (Lane). The trade pattern \emph{faʿʿāl} makes him the man of the door.
\end{cousinweb}

\begin{grammarlens}[title={What you've built}]
One more word for getting around.
\end{grammarlens}

\subsection*{Your turn}
\begin{itemize}
\raggedright
  \item \emph{Say it:} \emph{bawwāb}
  \item \emph{Say it:} \emph{bawwāb}, once more
  \item \emph{Say it:} \emph{ḥamala al-bawwāb amtiʿatanā ilā al-miṣʿad.} — The doorman carried our luggage to the lift.
  \item \emph{Recall:} say the Arabic for to set alight, then the Arabic for pine trees, then say \emph{bawwāb} again
\end{itemize}

\begin{tcolorbox}[breakable,colback=teal!4,colframe=teal!35!black,title={Before you move on}]
What does \ar{بواب} mean? (\textquotedblleft{}A doorman; a caretaker\textquotedblright{}.) Say it once more.
\end{tcolorbox}
\section[\texorpdfstring{\ar{ملاءة} (mulāʾa)}{ملاءة (mulāʾa)}]{\ar{ملاءة} (mulāʾa) — a bed sheet}

\label{lesson:AR-C443-mulaa}



\begin{quote}
Before the new one: say the Arabic for pine trees, then the Arabic for a doorman.
\end{quote}

\subsection*{You'll want to know: \ar{ملاءة}}
\textbf{\ar{ملاءة}} — \emph{mulāʾa} — \textquotedblleft{}a bed sheet\textquotedblright{}.

\begin{quote}
\ar{كانت} \ar{الملاءة} \ar{نظيفة} \ar{وناعمة}.
\end{quote}

\begin{quote}
\emph{kānat al-mulāʾa naẓīfa wa-nāʿima.}
\end{quote}

\begin{quote}
The sheet was clean and soft.
\end{quote}

\begin{cousinweb}
\textbf{\ar{ملاءة}} (\emph{mulāʾa}) is built on the native Arabic root \emph{m-l-ʾ}, \textquotedblleft{}to fill\textquotedblright{} (Lane). Lane gives \emph{mulāʾa} as a cloth wrapped about the body, often two pieces sewn together; modern Arabic uses it for a bed sheet.
\end{cousinweb}

\begin{grammarlens}[title={What you've built}]
One more word for getting around.
\end{grammarlens}

\subsection*{Your turn}
\begin{itemize}
\raggedright
  \item \emph{Say it:} \emph{mulāʾa}
  \item \emph{Say it:} \emph{mulāʾa}, once more
  \item \emph{Say it:} \emph{kānat al-mulāʾa naẓīfa wa-nāʿima.} — The sheet was clean and soft.
  \item \emph{Recall:} say the Arabic for pine trees, then the Arabic for a doorman, then say \emph{mulāʾa} again
\end{itemize}

\begin{tcolorbox}[breakable,colback=teal!4,colframe=teal!35!black,title={Before you move on}]
What does \ar{ملاءة} mean? (\textquotedblleft{}A bed sheet\textquotedblright{}.) Say it once more.
\end{tcolorbox}
\section[\texorpdfstring{\ar{نظافة} (naẓāfa)}{نظافة (naẓāfa)}]{\ar{نظافة} (naẓāfa) — cleanliness}

\label{lesson:AR-C443-nazafa}



\begin{quote}
Before the new one: say the Arabic for a doorman, then the Arabic for a bed sheet.
\end{quote}

\subsection*{You'll want to know: \ar{نظافة}}
\textbf{\ar{نظافة}} — \emph{naẓāfa} — \textquotedblleft{}cleanliness\textquotedblright{}.

\begin{quote}
\ar{كانت} \ar{نظافة} \ar{الغرفة} \ar{ممتازة}.
\end{quote}

\begin{quote}
\emph{kānat naẓāfat al-ghurfa mumtāza.}
\end{quote}

\begin{quote}
The cleanliness of the room was excellent.
\end{quote}

\begin{cousinweb}
\textbf{\ar{نظافة}} (\emph{naẓāfa}) is built on the native Arabic root \emph{n-ẓ-f}, \textquotedblleft{}to be clean\textquotedblright{} (Lane). It is the noun beside \emph{naẓīf}, \textquotedblleft{}clean\textquotedblright{}.
\end{cousinweb}

\begin{grammarlens}[title={What you've built}]
One more word for getting around.
\end{grammarlens}

\subsection*{Your turn}
\begin{itemize}
\raggedright
  \item \emph{Say it:} \emph{naẓāfa}
  \item \emph{Say it:} \emph{naẓāfa}, once more
  \item \emph{Say it:} \emph{kānat naẓāfat al-ghurfa mumtāza.} — The cleanliness of the room was excellent.
  \item \emph{Recall:} say the Arabic for a doorman, then the Arabic for a bed sheet, then say \emph{naẓāfa} again
\end{itemize}

\begin{tcolorbox}[breakable,colback=teal!4,colframe=teal!35!black,title={Before you move on}]
What does \ar{نظافة} mean? (\textquotedblleft{}Cleanliness\textquotedblright{}.) Say it once more.
\end{tcolorbox}
\section[\texorpdfstring{\ar{أطل} (aṭalla)}{أطل (aṭalla)}]{\ar{أطل} (aṭalla) — to overlook}

\label{lesson:AR-C443-atalla}



\begin{quote}
Before the new one: say the Arabic for a bed sheet, then the Arabic for cleanliness.
\end{quote}

\subsection*{You'll want to know: \ar{أطل}}
\textbf{\ar{أطل}} — \emph{aṭalla} — \textquotedblleft{}to overlook\textquotedblright{}.

\begin{quote}
\ar{تطل} \ar{غرفتنا} \ar{على} \ar{البحر}.
\end{quote}

\begin{quote}
\emph{tuṭillu ghurfatunā ʿalā al-baḥr.}
\end{quote}

\begin{quote}
Our room overlooks the sea.
\end{quote}

Present: \ar{يطل}, \emph{yuṭill}.

It takes \emph{ʿalā}: \emph{aṭalla ʿalā}, \textquotedblleft{}to look out over\textquotedblright{}.

\begin{cousinweb}
\textbf{\ar{أطل}} (\emph{aṭalla}) is usually traced to the native Arabic root \emph{ṭ-l-l}, \textquotedblleft{}to look down from above\textquotedblright{}. It is Form IV of a doubled root.
\end{cousinweb}

\begin{grammarlens}[title={What you've built}]
One more word for getting around.
\end{grammarlens}

\subsection*{Your turn}
\begin{itemize}
\raggedright
  \item \emph{Say it:} \emph{aṭalla}
  \item \emph{Say it:} \emph{aṭalla}, once more
  \item \emph{Say it:} \emph{tuṭillu ghurfatunā ʿalā al-baḥr.} — Our room overlooks the sea.
  \item \emph{Recall:} say the Arabic for a bed sheet, then the Arabic for cleanliness, then say \emph{aṭalla} again
\end{itemize}

\begin{tcolorbox}[breakable,colback=teal!4,colframe=teal!35!black,title={Before you move on}]
What does \ar{أطل} mean? (\textquotedblleft{}To overlook\textquotedblright{}.) Say it once more.
\end{tcolorbox}
\section[\texorpdfstring{\ar{خزنة} (khazna)}{خزنة (khazna)}]{\ar{خزنة} (khazna) — a strongbox; a safe}

\label{lesson:AR-C443-khazna}



\begin{quote}
Before the new one: say the Arabic for cleanliness, then the Arabic for to overlook.
\end{quote}

\subsection*{You'll want to know: \ar{خزنة}}
\textbf{\ar{خزنة}} — \emph{khazna} — \textquotedblleft{}a strongbox; a safe\textquotedblright{}.

\begin{quote}
\ar{وضعت} \ar{جوازي} \ar{في} \ar{خزنة} \ar{الغرفة}.
\end{quote}

\begin{quote}
\emph{waḍaʿtu jawāzī fī khaznat al-ghurfa.}
\end{quote}

\begin{quote}
I put my passport in the room safe.
\end{quote}

\begin{cousinweb}
\textbf{\ar{خزنة}} (\emph{khazna}) is built on the native Arabic root \emph{kh-z-n}, \textquotedblleft{}to store up\textquotedblright{} (Lane). It shares its root with \emph{khizāna}, \textquotedblleft{}a cupboard\textquotedblright{}.
\end{cousinweb}

\begin{grammarlens}[title={What you've built}]
One more word for getting around.
\end{grammarlens}

\subsection*{Your turn}
\begin{itemize}
\raggedright
  \item \emph{Say it:} \emph{khazna}
  \item \emph{Say it:} \emph{khazna}, once more
  \item \emph{Say it:} \emph{waḍaʿtu jawāzī fī khaznat al-ghurfa.} — I put my passport in the room safe.
  \item \emph{Recall:} say the Arabic for cleanliness, then the Arabic for to overlook, then say \emph{khazna} again
\end{itemize}

\begin{tcolorbox}[breakable,colback=teal!4,colframe=teal!35!black,title={Before you move on}]
What does \ar{خزنة} mean? (\textquotedblleft{}A strongbox; a safe\textquotedblright{}.) Say it once more.
\end{tcolorbox}
